\section{Potential Defenses}
In this section, we discuss whether our attack is resistant to potential defenses.

%we discuss and evaluate potential methods to defend against our attack.
\begin{figure}[t]
\centering
\includegraphics[width=0.4\textwidth]{fig/dp.pdf}
%\vspace{-0.5em}
\caption{Token accuracy and perplexity under defense based on activation obfuscation on Skytrax~\cite{skytrax} dataset. }
%The X-axis represents different Gaussian noise.}
\label{fig:DP}
\vspace{-1em}
\end{figure}



\noindent
\textbf{Activation obfuscation. }
Adding random noise to obfuscate sensitive information is a general privacy protection paradigm. 
Without loss of generality, we consider a scheme where each participant adds random Gaussian noise $\epsilon$ to the activation before sending it. We refer to this scheme as \emph{activation obfuscation}. 
We utilize \textit{perplexity}~\cite{perplexity}, measured on a Llama-7B model, to assess the quality of the text generated by the activation obfuscation scheme.  From Figure~\ref{fig:DP} , we can see that when the variance $\sigma$ of noise  $\epsilon$  is small (e.g., 0.1), PIA still recovers most of the prompt information (token accuracy 87.66\%). When noise is greater, the inversion accuracy decreases slightly, while the quality of the generated text is severely reduced. When $\sigma=1$, the perplexity increases to 14.2. These results indicate that \emph{activation obfuscation}  does not achieve a satisfactory trade-off between privacy protection and LLM utility. 

%There is an inherent trade-off here: increasing noise improves privacy but reduces the quality of the generated text.
%activation obfuscation does not seem to be practical in defending our attack. 

\noindent
\textbf{Activation quantization.}
\emph{Activation quantization} is a widely used technique for minimizing communication in collaborative inference and fine-tuning systems~\cite{pact, cocktail}. 
\revise{Unlike conventional LLM quantization, in activation quantization, the user does not finetune the model weights, as they cannot modify the portions of the weights controlled by other users.  We consider it as a potential defense targetting to reduce information leakage to attackers.} We conduct experiments in the setting where each participant quantizes their activation to 8 bits and 4 bits before sending it. \emph{Activation quantization}  also involves a trade-off between text quality and privacy. From Table~\ref{tab:quantization} in Appendix, we observe that inversion still achieves high accuracy under 8-bit quantization. For 4-bit quantization, the inversion accuracy drops significantly, but the text quality is also severely damaged (perplexity 28.3). These results indicate that \emph{activation quantization} is not an effective defense against our attack.

%More details and our experiments can be found in Appendix~\ref{appendix:defense}.

%Therefore, we utilize \textbf{Perplexity}~\cite{perplexity}, measured on the Vicuna-13B model, to assess the quality of the text generated by 
%We utilize \textbf{Perplexity}~\cite{perplexity}, measured on the Vicuna-13B model, to assess the text quality. This metric is commonly employed in NLP domain. Smaller perplexity usually indicates a better quality. 


\noindent
\textbf{Homomorphic encryption.} Leveraging homomorphic encryption~\cite{HE} enables participants to perform DNN computations on encrypted inputs, thereby preserving privacy. However, the computational overhead associated with homomorphically computing an LLM is prohibitively high~\cite{chen-etal-2022-x}, making real-world deployment challenging.




%is effective at secure remote computation. In our scenario, each participant could deploy a TEE to compute the inference. Specifically, a TLS-encrypted channel between enclaves can be established through remote attestation~\cite{knauth2018integrating}. Once established, the activations are encrypted against the adversary outside the enclave during transmission. However, the enclave can still use the activation for subsequent computations, ensuring security. This solution could provide privacy protection for the client prompt. However, it requires special architectural support, which prohibits its application to many participants. Another severe problem is that TEEs introduce considerable overhead in serving LLMs~\cite{mai2023honeycomb}, making decentralized TEE-based inference impractical. 


%\noindent
%\textbf{Homomorphic Encryption. }
% allows participants to perform DNN computations directly on encrypted inputs. The resulting activations are also encrypted, thus preventing privacy leakage. A drawback of homomorphic encryption is  huge inefficiency and large memory overhead. The overhead of homomorphically computing an LLM is too high~\cite{chen-etal-2022-x} , making it hard to deploy in the real world.